% 张三慧《大学物理学》第4版《力学、热学》第5章（扫描PDF物理页185--189）
% 章末习题5.1--5.28；按原题号保留数据、图示关系、分问。
% 需要原始图5.30--5.39验证的条目继续列为待审候选。
\section*{第5章\quad 刚体的转动}
\subsection*{章末习题5.1--5.28}

\paragraph{习题 5.1}
掷铁饼运动员手持铁饼转动 \(1.25\) 圈后松手，
此刻铁饼速率达到 \(v=25\,\mathrm{m/s}\)。
设转动时铁饼沿半径为 \(R=1.0\,\mathrm m\) 的圆周运动
并且均匀加速，求：
（1）铁饼离手时的角速度；
（2）铁饼的角加速度；
（3）铁饼在手中加速的时间（把铁饼视为质点）。
\textit{条件说明：}原题未明说起始角速度；求后两问需采用
从静止开始转动的通常模型，否则答案不唯一。

\paragraph{习题 5.2}
一汽车发动机的主轴的转速在 \(7.0\,\mathrm s\) 内
由 \(200\,\mathrm{r/min}\) 均匀地增加到
\(3000\,\mathrm{r/min}\)。求：
（1）在这段时间内主轴的初角速度和末角速度以及角加速度；
（2）这段时间内主轴转过的角度和圈数。

\paragraph{习题 5.3}
地球自转是逐渐变慢的。
在1987年完成365次自转比1900年长 \(1.14\,\mathrm s\)。
求在1900年到1987年这段时间内，
地球自转的平均角加速度。

\paragraph{习题 5.4}
求位于北纬 \(40^\circ\) 的颐和园排云殿
（以图5.30中的 \(P\) 点表示）
相对于地心参考系的线速度与加速度的数值与方向。
\textit{图示：}地球近似绕南北自转轴以角速度 \(\omega\) 转动，
\(P\) 到地球自转轴的距离为 \(R\cos40^\circ\)。

\paragraph{习题 5.5}
水分子的形状如图5.31所示。
从光谱分析得出其对 \(AA'\) 轴的转动惯量是
\(J_{AA'}=1.93\times10^{-47}\,\mathrm{kg\,m^2}\)，
对 \(BB'\) 轴的转动惯量是
\(J_{BB'}=1.14\times10^{-47}\,\mathrm{kg\,m^2}\)。
试由此数据和各原子的质量求出氧原子与氢原子间的距离
\(d\) 和夹角 \(\theta\)。
假设各原子都可当质点处理。
\textit{图示：}氧原子 \(O\) 位于两对称轴的交点，
两根 \(OH\) 键长均为 \(d\)，
各与水平 \(AA'\) 正轴成夹角 \(\pm\theta\)。

\paragraph{习题 5.6}
\(C_{60}\)（Fullerene，富勒烯）分子由60个碳原子组成，
这些碳原子各位于一个球形32面体的60个顶角上（图5.32），
此球体的直径为 \(71\,\mathrm{nm}\)。
（1）按均匀球面计算，此球形分子对其一个直径的转动惯量是多少？
（2）在室温下一个 \(C_{60}\) 分子的自转动能为
\(6.21\times10^{-21}\,\mathrm J\)。求它的自转频率。
\textit{原文疑点：}扫描页185--186明确印为「71 nm」，
与常见 \(C_{60}\) 分子直径约 \(0.71\,\mathrm{nm}\)
相差100倍；本题保留原印刷数据，分别对照计算并待复核，
绝不静默改写原题。

\paragraph{习题 5.7}
一个氧原子的质量是 \(2.66\times10^{-26}\,\mathrm{kg}\)，
一个氧分子中两个氧原子的中心相距
\(1.21\times10^{-10}\,\mathrm m\)。
求氧分子相对于通过其质心并垂直于二原子连线的轴的转动惯量。
如果一个氧分子相对于此轴的转动动能是
\(2.06\times10^{-21}\,\mathrm J\)，
它绕此轴的转动周期是多少？

\paragraph{习题 5.8}
一个哑铃由两个质量为 \(m\)、半径为 \(R\) 的铁球
和中间一根长 \(\ell\) 的连杆组成（图5.33）。
和铁球的质量相比，连杆的质量可以忽略。
求此哑铃对于通过连杆中心并和它垂直的轴的转动惯量。
它对于通过两球连心线的轴的转动惯量又是多少？
\textit{图示：}杆长 \(\ell\) 是两球表面之间的距离，
球心间距离为 \(\ell+2R\)。

\paragraph{习题 5.9}
在伦敦的英国议会塔楼上的大本钟的分针长
\(4.50\,\mathrm m\)，质量为 \(100\,\mathrm{kg}\)；
时针长 \(2.70\,\mathrm m\)，质量为 \(60.0\,\mathrm{kg}\)。
二者对中心轴的角动量和转动动能各是多少？
将二者都当成均匀直棒处理。

\paragraph{习题 5.10}
从一个半径为 \(R\) 的均匀薄板上挖去
一个直径为 \(R\) 的圆板，
所形成的圆洞中心在距原薄板中心 \(R/2\) 处
（图5.34），剩余薄板的质量为 \(m\)。
求此时薄板对于通过原中心而与板面垂直的轴的转动惯量。

\paragraph{习题 5.11}
如图5.35所示，两物体质量分别为 \(m_1,m_2\)，
定滑轮质量为 \(m\)、半径为 \(r\)，可视为均匀圆盘。
已知 \(m_2\) 与桌面间的滑动摩擦系数为 \(\mu_k\)，
求 \(m_1\) 下落的加速度和两段绳子中的张力各是多少？
设绳子和滑轮间无相对滑动，滑轮轴承的摩擦力忽略不计。

\paragraph{习题 5.12}
一根均匀米尺，在 \(60\,\mathrm{cm}\) 刻度处被钉到墙上，
且可以在竖直平面内自由转动。
先用手使米尺保持水平，然后释放。
求刚释放时米尺的角加速度和米尺到竖直位置时的角速度各是多少？

\paragraph{习题 5.13}
从质元的动能表示式
\(\Delta E_k=\tfrac12\Delta m\,v^2\) 出发，
导出刚体绕定轴转动的动能表示式
\(E_k=\tfrac12J\omega^2\)。

\paragraph{习题 5.14}
唱机的转盘绕着通过盘心的固定竖直轴转动，
唱片放上去后将受转盘的摩擦力作用而随转盘转动（图5.36）。
设唱片可以看成是半径为 \(R\) 的均匀圆盘，质量为 \(m\)，
唱片和转盘之间的滑动摩擦系数为 \(\mu_k\)。
转盘原来以角速度 \(\omega\) 匀速转动，唱片刚放上去时是静止的。
唱片达到角速度 \(\omega\) 需要多长时间？
在这段时间内，转盘保持角速度 \(\omega\) 不变，
驱动它的电动机做了多少功？
唱片获得了多大动能？

\paragraph{习题 5.15}
坐在转椅上的人手握哑铃（图5.12）。
两臂伸直时，人、哑铃和转椅系统对竖直轴的转动惯量为
\(J_1=2\,\mathrm{kg\,m^2}\)。
在外人推动后，此系统开始以 \(n_1=15\,\mathrm{r/min}\) 转动。
当人的两臂收回，使系统的转动惯量变为
\(J_2=0.80\,\mathrm{kg\,m^2}\) 时，
它的转速 \(n_2\) 是多大？
两臂收回过程中，系统的机械能是否守恒？
什么力做了功？做功多少？
设轴上摩擦忽略不计。

\paragraph{习题 5.16}
图5.37中均匀杆长 \(L=0.40\,\mathrm m\)，
质量 \(M=1.0\,\mathrm{kg}\)，
由其上端的光滑水平轴吊起而处于静止。
今有一质量 \(m=8.0\,\mathrm g\) 的子弹
以 \(v_0=200\,\mathrm{m/s}\) 的速率
水平射入杆中而不复射出，
射入点在轴下 \(d=3L/4\) 处。
求：（1）子弹停在杆中时杆的角速度；
（2）杆的最大偏转角。

\paragraph{习题 5.17}
一转台绕竖直固定轴转动，
每转一周所需时间为 \(t=10\,\mathrm s\)，
转台对轴的转动惯量为 \(J=1200\,\mathrm{kg\,m^2}\)。
一质量为 \(M=80\,\mathrm{kg}\) 的人开始时站在转台中心，
随后沿半径向外跑去，
当人离转台中心 \(r=2\,\mathrm m\) 时
转台的角速度是多大？

\paragraph{习题 5.18}
两辆质量都是 \(1200\,\mathrm{kg}\) 的汽车
在平直公路上都以 \(72\,\mathrm{km/h}\)
的高速迎面开行。
由于两车质心轨道间距离太小，
仅为 \(0.5\,\mathrm m\)，因而发生碰撞。
碰后两车扣在一起，此残体对于其质心的转动惯量为
\(2500\,\mathrm{kg\,m^2}\)，求：
（1）两车扣在一起时的旋转角速度；
（2）由于碰撞而损失的机械能。

\paragraph{习题 5.19}
宇宙飞船中有三个宇航员穿着船舱环形内壁
按同一方向跑动以产生人造重力。
（1）如果想使人造重力等于他们在地面上时受的自然重力，
那么他们跑动的速率应多大？
设他们的质心运动的半径为 \(2.5\,\mathrm m\)，
人体当质点处理。
（2）如果飞船最初未动，当宇航员按上面速率跑动时，
飞船将以多大角速度旋转？
设每个宇航员的质量为 \(70\,\mathrm{kg}\)，
飞船船体对于其纵轴的转动惯量为
\(3\times10^5\,\mathrm{kg\,m^2}\)。
（3）要使飞船转过 \(30^\circ\)，宇航员需要跑几圈？
\textit{参考系说明：}跑速是相对船舱还是惯性系的速度
会造成小幅差异，必须注明所用模型。

\paragraph{习题 5.20}
把太阳当成均匀球体，
试由本书的「数值表」给出的有关数据
计算太阳的角动量。
太阳的角动量大约是太阳系总角动量
（\(3.3\times10^{43}\,\mathrm{kg\,m^2/s}\)）的百分之几？

\paragraph{习题 5.21}
蟹状星云（图5.38）中心是一颗脉冲星，
代号PSR 0531+21。
它以十分确定的周期（\(0.033\,\mathrm s\)）
向地球发射电磁波脉冲，
由中子密聚而成，脉冲周期就是它的转动周期。
实测发现，上述中子星的周期以
\(1.26\times10^{-5}\,\mathrm{s/a}\) 的速率增大。
（1）求此中子星的自转角加速度；
（2）设此中子星质量为 \(1.5\times10^{30}\,\mathrm{kg}\)
（近似太阳的质量）、半径为 \(10\,\mathrm{km}\)，
求它的转动动能以多大的速率（以 \(\mathrm{J/s}\) 计）减小。
（这减少的转动动能就转变为蟹状星云向外辐射的能量。）
（3）若这一能量变化率保持不变，
该中子星经过多长时间将停止转动？
设中子星可作均匀球体处理。

\paragraph{习题 5.22}
地球对自转轴的转动惯量是 \(0.33MR^2\)，
其中 \(M\) 是地球的质量，\(R\) 是地球的半径。
求地球的自转动能。
由于潮汐对海岸的摩擦作用，
地球自转的速度逐渐减小，
每百万年自转周期增加 \(16\,\mathrm s\)。
这样，地球自转动能的减小相当于摩擦消耗多大的功率？
一年内消耗的能量相当于我国2004年发电量的几倍？
潮汐对地球的平均力矩多大？

\paragraph{习题 5.23}
太阳的热核燃料耗尽时，它将急速塌缩成半径等于地球半径的一颗白矮星。
如果不计质量散失，那时太阳的转动周期将变为多少？
太阳和白矮星均按均匀球体计算，
目前太阳的自转周期按 \(26\,\mathrm d\) 计。

\paragraph{习题 5.24}
证明：圆盘在平面上无滑动地滚动时，
其上各点相对于平面的速度 \(\boldsymbol v\)
和圆盘的转动角速度 \(\boldsymbol\omega\) 有
\[
 \boldsymbol v=\boldsymbol\omega\times\boldsymbol r_P,
\]
其中 \(\boldsymbol\omega\) 垂直于盘面，
其方向按右手螺旋定则决定，
\(\boldsymbol r_P\) 为从圆盘与平面瞬时接触点 \(P\)
到各点的径矢。
故接触点 \(P\) 是圆盘的瞬时转动中心。

\paragraph{习题 5.25}
绕有电缆的大轮轴总质量
\(M=1000\,\mathrm{kg}\)，
轮半径 \(R_1=1.00\,\mathrm m\)，
电缆在轴上绕至半径为
\(R_2=0.60\,\mathrm m\) 处。
设此时整个轮轴对其中心轴的转动惯量
\(J_C=300\,\mathrm{kg\,m^2}\)。
今用 \(F=2000\,\mathrm N\) 的力
在底部沿水平方向拉电缆，
如果轮轴在路面上无滑动地滚动，
它将向哪个方向滚动？
滚动的角加速度多大？
质心前进的加速度多大？
轮轴受到地面的摩擦力多大？

\paragraph{习题 5.26}
小学生爱玩的「悠悠球」是把一条绕在
一个扁圆柱体的圆柱面上的线绳，用手抓住线绳的一端
上下抖动，使扁圆柱体上下运动的同时还不停地绕其水平轴转动。
下面为简单起见，假定线就绕在扁圆柱体的圆柱面上并沿圆柱体的质量为
\(m\)、半径为 \(R\)。
（1）若手不动，让圆柱体沿竖直的线自行滚下，
它下降的加速度多大？手需用多大的力提住线端？
（2）若要使圆柱体停留在一定高度上，
手需用多大的力向上提起线端？
圆柱体转动的角加速度多大？
（3）手若用 \(2mg\) 的力向上提起线端，
则圆柱体上升的加速度多大？
手向上提起的加速度多大？
\textit{原题文字和几何：}线沿圆柱体半径 \(R\) 绕卷，
质量均匀，需区分圆柱体的质心加速度与线端加速度。

\paragraph{习题 5.27}
地球的自转轴与它绕太阳的轨道平面的垂线间的夹角是
\(23.5^\circ\)（图5.39）。
由于太阳和月亮对地球的引力产生力矩，
地球的自转轴绕轨道平面的垂线进动，
进动一周需时间约 \(26000\,\mathrm a\)。
已知地球绕自转轴的转动惯量为
\(J=8.05\times10^{37}\,\mathrm{kg\,m^2}\)。
求地球自旋角动量矢量变化率的大小，
即 \(|d\boldsymbol L/dt|\)，
并求太阳和月亮对地球的合力矩多大。

\paragraph{习题 5.28}
一般船中装的回转稳定器是一个质量为
\(50.0\,\mathrm t\)、半径为 \(2.00\,\mathrm m\) 的固体圆盘。
它绕着一个竖直轴以 \(800\,\mathrm{r/min}\)
的角速度转动。
（1）如果用恒定输入功率
\(7.46\times10^4\,\mathrm W\) 启动，
要经过多少时间才能使它从静止达到上述额定转速？
（2）如果船的轴在船的纵向竖直平面内以
\(1.00^\circ/\mathrm s\) 的角速度进动，
说明船体受到了左倾或右倾的多大力矩？
\textit{方向说明：}未给定转盘自旋正方向与船体进动方向的箭头，
能唯一确定力矩大小，左右倾方向需结合原图的符号约定。
